% id: 105-09
% book: 베이직쎈 미적분1 · 06 도함수의 활용 (3)
% section: 학교 시험 기출로 실전 감각 UP!
% figure: tikz:fig-105-09
% answer: 
% answer_source: 
% variant_level: 0 (원본 전사)
% 이 파일은 build.mjs 가 items.json 에서 만든다. 고칠 때는 items.json 을 고친다.
\smash{\raisebox{1.5mm}{\dmkichul{베이직쎈 미적분1 105쪽 09번}}}
\noindent\begin{minipage}[t]{0.58\linewidth}\vspace{0pt}\raggedright 원점을 출발하여 수직선 위를 움직이는 점~$\pt{P}$의 시각 $t$에서의 위치 $x(t)$의 그래프가 오른쪽 그림과 같을 때, 다음 중 옳지 \underline{않은} 것은?\end{minipage}\hfill\begin{minipage}[t]{0.4\linewidth}\vspace{0pt}\centering \resizebox{\ifdim\width>\linewidth\linewidth\else\width\fi}{!}{% 105-09(105쪽) — 원점을 출발한 점 P 의 시각 t 에서의 위치 x(t) 의 그래프(t=b 에서 극소, t=c 에서 0, t=d 에서 극대, t=e 에서 t축에 닿았다가 다시 증가). 스캔 화질이 낮아 TikZ 로 다시 그림.
% 원문에 있는 것만 옮긴다: 곡선 · 안내 점선 세 줄(a, b, d) · 라벨 x(t), t, O, a, b, c, d, e. 원문에 눈금값이 없으므로 적지 않는다.
\begin{tikzpicture}[x=0.32cm, y=0.32cm, line width=0.6pt]
  \draw[->, >=stealth, line width=0.7pt] (-0.75,0) -- (7.75,0) node[below=1pt, font=\small] {$t$};
  \draw[->, >=stealth, line width=0.7pt] (0,-1.75) -- (0,3.05) node[left=1pt, font=\small] {$x(t)$};
  \draw[densely dashed, line width=0.4pt] (1.05,0) -- (1.05,-1.12);
  \draw[densely dashed, line width=0.4pt] (1.60,0) -- (1.60,-1.25);
  \draw[densely dashed, line width=0.4pt] (4.25,0) -- (4.25,1.30);
  \draw plot[smooth, tension=0.7] coordinates {(0,0) (0.45,-0.55) (0.80,-0.94) (1.20,-1.19) (1.60,-1.25) (2.05,-1.15) (2.55,-0.85) (3.15,0) (3.60,0.62) (4.00,1.13) (4.25,1.30) (4.65,1.18) (5.05,0.76) (5.32,0.34) (5.48,0.11) (5.60,0) (5.74,0.08) (5.95,0.32) (6.25,0.90) (6.70,1.88) (7.00,2.70)};
  \node[below left=-2pt and -2.5pt, font=\small] at (0,0) {$\pt{O}$};
  \node[above=0.5pt, font=\small] at (1.05,0) {$a$};
  \node[above=0.5pt, font=\small] at (1.60,0) {$b$};
  \node[below=1pt, font=\small] at (3.15,0) {$c$};
  \node[below=1pt, font=\small] at (4.25,0) {$d$};
  \node[below=1pt, font=\small] at (5.60,0) {$e$};
\end{tikzpicture}
}\end{minipage}\par
\par\medskip\par\noindent\hangindent=1.4em\hangafter=1 {\small ①}\ $t=b$에서 점~$\pt{P}$의 속도는 $0$이다.\par\noindent\hangindent=1.4em\hangafter=1 {\small ②}\ $t=c$에서 점~$\pt{P}$는 원점을 지난다.\par\noindent\hangindent=1.4em\hangafter=1 {\small ③}\ $d<t<e$에서 점~$\pt{P}$는 음의 방향으로 움직인다.\par\noindent\hangindent=1.4em\hangafter=1 {\small ④}\ $0<t<e$에서 $t=d$일 때 점~$\pt{P}$의 속도가 최대이다.\par\noindent\hangindent=1.4em\hangafter=1 {\small ⑤}\ $0<t<e$에서 점~$\pt{P}$는 운동 방향을 두 번 바꾼다.\par\medskip
